\documentclass[11pt]{uebungsblatt}

\blatttitel{Bonusübungsblatt -- Zur Selbstüberprüfung (unbenotet)}
\blattuntertitel{\LaTeX }
\excludecomment{loesung}

\usepackage{graphicx}
\usepackage{xcolor}

\def\url#1{\expandafter \textcolor{blue}{#1} \csname #1\endcsname}

\header{\textbf{Algorithmen und Datenstrukturen}\hfill\textbf{Wintersemester 2026/27}\\
	Prof.~Dr.~Florin~Manea \hfill Georg-August-Universität Göttingen \\
	M.Sc.~Timo~Specht \hfill Institut für Informatik\\
	\strut\hfill Lehrstuhl für Theoretische Informatik\\
	\rule{\textwidth}{0.1mm}}
\introtext{%
	
}

\begin{document}

\begin{aufgabe}
	Eine kleine Einführung zur Verwendung von \LaTeX\ ist für Sie im Stud.IP unter Medien hochgeladen.

	In der Stud.IP Veranstaltung zur Vorlesung befindet sich unter \textit{Dateien} $\rightarrow$ \textit{Übung} $\rightarrow$ \textit{Latex} $\rightarrow$ \textit{template.zip} eine allgemeine Vorlage für die Anfertigung Ihrer Abgaben.
	Laden Sie das Archiv herunter und entpacken Sie es.
	Bei den Dateien \textit{uebungsblatt.sty} und \textit{uebungsblatt.cls} handelt es sich um Klassen- bzw. Style-Files,
	die Sie zum kompilieren benötigen, aber \textbf{nicht} weiter bearbeiten sollen.
	Die Datei \textit{uebung.tex} soll von Ihnen angepasst werden,
	indem Sie im Header an den entsprechenden Stellen {Student:in1} und {Student:in2} durch die Namen der abgebenden Studierenden ersetzen.
	Ergänzen Sie außerdem die entsprechende Blattnummer.

	Zum Editieren und Kompilieren können Sie beispielsweise den \LaTeX-Editor des GWDG verwenden.
	Dieser ist unter \url{https://sharelatex.gwdg.de/} zu finden.
	Hier müssen Sie lediglich Ihre Datei zusammen mit den beiden Klassen-/Style-Files in einem neuen Projekt hochladen.

	Für zukünftige Abgaben wird stets ein entsprechendes Template mit den aktuellen Aufgaben im Latex-Ordner hochgeladen werden,
	welches Sie nutzen können.
\end{aufgabe}

\textbf{Verwenden Sie Ihre angepasste Vorlage aus Aufgabe 1 nun für die folgenden Aufgaben.}

\begin{aufgabe}
	Setzen Sie in Ihrer \LaTeX-Datei folgenden Text inklusive Zeichenformatierung und der Tabelle. Verwenden Sie dabei mindestens an einer Stelle eine Gruppe.

	\bigskip
	Um einen \LaTeX-Text zu setzen, kann man sich vieler Mittel bedienen.
	\textit{Sehen Sie sich zum Beispiel diesen Satz an, er ist kursiv!}
	\textsf{\textbf{Aber auch serifenlos und fett können einzelne Textelemente gesetzt werden.}}

	Absätze innerhalb eines Texts können wir einfach durch eine Leerzeile erreichen.
	Darüber hinaus werden Zeilenumbrüche innerhalb eines \LaTeX-Textes einfach ignoriert.
	Zeichen wie \& oder \textbackslash\ müssen normalerweise speziell gesetzt werden,
	damit diese in einem Text verwendet werden können.

	\begingroup
	\LARGE
	Um Text besonders hervorzuheben, können wir die Schriftgröße auf LARGE setzen.
	\endgroup
	Dabei ist allerdings wichtig, dass wir die Schriftgröße für darauffolgenden Text auch wieder auf ein normales Maß runterstellen.
	Um dies zu erreichen können wir beispielsweise mit Gruppen arbeiten.

	Zum Abschluss folgt hier nun eine Tabelle mit einer kleinen Übersicht möglicher Zeichensätze.

	\begin{table}[H]
		\centering
		\begin{tabular}{lll}
			\multicolumn{2}{r}{\textbf{Familie}}                                                              \\
			\texttt{\string\rmfamily} & \texttt{\string\textrm\{\}} & \textrm{normal}                         \\
			\texttt{\string\sffamily} & \texttt{\string\textsf\{\}} & \textsf{serifenfrei}                    \\
			\texttt{\string\ttfamily} & \texttt{\string\texttt\{\}} & \texttt{Schreibmaschine} \vspace{0.5cm} \\
			\multicolumn{2}{r}{\textbf{Stärke}}                                                               \\
			\verb|\mdseries|          & \verb|\textmd{}|            & \textmd{normal}                         \\
			\verb|\bfseries|          & \verb|\textbf{}|            & \textbf{fett} \vspace{0.5cm}            \\
			\multicolumn{2}{r}{\textbf{Varianten}}                                                            \\
			\verb|\upshape|           & \verb|\textup{}|            & \textup{aufrecht}                       \\
			\verb|\itshape|           & \verb|\textit{}|            & \textit{kursiv}                         \\
			\verb|\scshape|           & \verb|\textsc{}|            & \textsc{Kapitälchen}                    \\
		\end{tabular}
	\end{table}
\end{aufgabe}

\newpage
\begin{aufgabe}
	Setzen Sie in Ihrer \LaTeX-Datei die folgenden mathematischen Ausdrücke:

	\begin{enumerate}[a)]
		\item $$ f \in \mathcal{O}(n), \quad g \in \Omega(n \log n) $$
		\item $$ \forall k \in \mathbb{N} \colon s \sim_k t \vcentcolon\Leftrightarrow \mathcal{SF}_k(s) = \mathcal{SF}_k(t) $$
		\item $$ f \colon \mathbb{Z} \rightarrow \mathbb{R}, \quad f(x) \coloneqq \sqrt[5]{\lvert x - 5 \rvert} $$
		\item $$ \underbrace{\varDelta \frac{\gamma_1^3}{2 m_0} = \sum_i^x \frac{i}{2\varepsilon}}_{\text{\Large{This is some nice equation stuff!}}} $$
		\item $$ fib(n) =
			      \begin{cases}
				      1                   & \text{for } n \in \{1,2\} \\
				      fib(n-1) + fib(n-2) & \text{for } n \geq 3
			      \end{cases}
		      $$
	\end{enumerate}
\end{aufgabe}

\begin{aufgabe}
	Fügen Sie für diese Aufgabe ein beliebiges Bild
	(z.B. ein Gruß, ein lustiges Meme, süße Katzenbabys, ein nettes Selfie, was auch immer)
	in Ihre \LaTeX-Datei ein.
	Wählen Sie eine geeignete Bildgröße.
	Versehen Sie das Bild mit einem Untertitel (auch caption genannt).

	Hinweis: Das Bild muss als pdf, jpg oder png vorliegen, darf maximal 1 MB groß sein und muss sich im selben Verzeichnis wie Ihre \LaTeX-Datei befinden.
\end{aufgabe}

\newpage
\begin{loesung}
	\begin{figure}[h]
		\centering
		\includegraphics[width=0.5\linewidth]{photo.png}
		\caption{Ungewöhnliche Freundschaft}
	\end{figure}
\end{loesung}

\newpage
%=====================================================================
\begin{cut}


	\begin{aufgabe}
		\begin{loesung}


		\end{loesung}
	\end{aufgabe}


\end{cut}
%=====================================================================

\end{document}